\section{Related Work}

\subsection{Association and reactivation in online trackers}

Tracking-by-detection pipelines pair a detector with a data association
rule. SORT \cite{bewley2016sort} fixed the modern form: a Kalman filter
predicts each track, and a Hungarian assignment on box overlap consumes
the detections. DeepSORT \cite{wojke2017deepsort} added an appearance
metric to the same cascade. ByteTrack \cite{zhang2022bytetrack} broke the
rule that low-confidence detections are discarded, and it introduced the
track buffer that makes reactivation routine: a track survives thirty
frames of absence and any later detection above the association threshold
is re-attached. BoT-SORT \cite{aharon2022botsort} improved the motion model
and compensated camera movement; OC-SORT \cite{cao2023ocsort} replaced the
filter state used in matching with observation-centered estimates.

These systems change what counts as a good match. Reconnection is a
by-product of that choice: the tracker keeps a lost track alive, and
whatever clears the threshold is accepted. The buffer length and the
threshold decide how often this happens, and neither is a statement about
whether a particular return is the same person. That is the gap we
address.

In a two-stage cascade the reactivation decision is made by the same scalar
threshold that governs ordinary association. That threshold is tuned for
the bulk of assignments, most of which involve no gap at all, so a return
after eight frames and a return after one frame are judged by the same
number; the only per-event variation comes from whatever motion cost the
filter happened to predict. Nothing in the cascade asks whether the
evidence at this particular proposal is strong. A learned per-proposal
score asks exactly that.

\subsection{Appearance-based re-identification}

The standard answer to an ambiguous return is an appearance descriptor.
DeepSORT \cite{wojke2017deepsort} introduced the idea, and later trackers
keep a re-identification branch alongside motion, or link tracklets after
tracking with a learned similarity. Where a trained embedding model and
enough pixels per crop are available, this signal is strong, and it
addresses occlusion gaps that geometry cannot resolve.

A descriptor is not free of failure modes of its own. Clothing changes,
matching uniforms and crowds of near-identical people all degrade it, and a
gallery has to be maintained across the sequence rather than read once. Its
cost also scales with the number of active tracks instead of the number of
reactivations, which matters in long sequences where most tracks live for a
few seconds.

It also imposes conditions we do not assume. It requires a network, a
feature gallery or a re-identification budget, and crops of usable
resolution at both ends of the gap. On sparse annotation grids and on
small or heavily occluded boxes those conditions often fail together. We
therefore restrict the study to detection-only evidence and treat
appearance as the natural extension rather than a missing component. The
comparison we do make is against an association score from a published
tracker, BoostTrack \cite{stanjoevic2024boosttrack}, ported into the same
host so that only the scoring component changes.

\subsection{Selection and risk control}

Gating proposals is a selection problem: a threshold is chosen from data,
and the chosen threshold's measured risk is optimistic by construction.
Learn-then-test \cite{angelopoulos2021learn} calibrates this choice over a
family of policies and returns a valid risk level. Closely related work
bounds the coverage of candidate transition sets in tracking
\cite{li2022pac} and propagates detector-box uncertainty through the
tracking pipeline \cite{su2024motcup}. We use the vocabulary of that line
of work to describe
what our selection rule does, and we report the selected policies'
observed failure proportions without claiming a calibrated bound. What
distinguishes the present study is not the statistics but the object being
measured: committed reactivations after closed-loop replay, counted
separately from the gap recoveries they produce.

Selection and filtering are easy to conflate, and the difference decides
what a fair comparison looks like. A threshold selects one operating point
for every proposal; a score decides each proposal on its own evidence. Two
rules that admit the same total number of actions can still commit
different actions, so a comparison that matches only the totals still
leaves the per-proposal decision unexamined. That is why we match admission
counts source by source and compare the paired decisions rather than the
aggregate rate alone.

\subsection{Benchmarks and how tracking is scored}

Three benchmarks shape what a tracking result is allowed to claim. MOT16
\cite{milan2016mot16} standardised annotation and made crowded pedestrian
sequences the default test, where clothing and build vary enough that an
appearance model has something to learn. DanceTrack
\cite{sun2022dancetrack} removes that advantage by design: its dancers wear
near-identical costumes, so appearance carries almost no information and
association has to rest on motion. PersonPath22
\cite{shuai2022personpath22}, the dataset used here, sits elsewhere again,
with sparse manual keyframes, a frozen detector and long sequences in which
the same person leaves the frame and comes back.

Metrics moved with the benchmarks. HOTA \cite{luiten2021hota} separates
detection accuracy from association accuracy and reports both before
combining them, which matters when a method changes only the association
step. PersonPath22 supports a measurement that neither benchmark offers
directly: because each keyframe carries both a box and an identity, a
detector gap can be located exactly, and its recovery counted separately
from the association quality of the sequence as a whole. That separation is
what lets this paper report an action-level failure proportion next to a
trajectory-level score instead of choosing between them.
